\documentclass[11pt]{article}
\usepackage[localtheorems]{raeez-math-template}
\newtheorem{theorem}{Theorem}[section]
\newtheorem{proposition}[theorem]{Proposition}
\newtheorem{lemma}[theorem]{Lemma}
\newtheorem{definition}[theorem]{Definition}
\newtheorem{remark}[theorem]{Remark}
\newtheorem{corollary}[theorem]{Corollary}
\newtheorem*{attack}{Attack}
\newtheorem*{heal}{Heal}

\title{BKM Serre Vanishing at Depth Three and Root-of-Unity Degeneracy}
\author{Raeez Lorgat}
\date{2026-04-22}

\begin{document}
\maketitle

\begin{abstract}
Two claims sit at the frontier of the Calabi--Yau-to-chiral functor
$\Phi\colon \mathrm{CY}\text{-cat}_d \to \mathrm{ChirAlg}$ at
$d = 3$. First: the vanishing of the second Serre polynomial
$P_2(D)$ for the Borcherds--Kac--Moody (BKM) superalgebra
$\mathfrak{g}_{\Delta_5}$ with Gritsenko denominator $\Delta_5$, at
depth $D = 3$ in the imaginary-root stratification. Second: the
count of $324$ finite-dimensional irreducibles of the quantum
toroidal algebra $U_{q,t}(\widehat{\widehat{\mathfrak{gl}}}_1)$ at
the second root of unity $q = -1$, and the degeneracy of the
$S$-matrix on its abelian sector. We subject both to five
attack--heal cycles and identify, behind each, a precise statement
about $6$d holomorphic Chern--Simons boundary conditions at
$\Omega$-background locus $\varepsilon_1 \varepsilon_2 = 0$.
\end{abstract}

\tableofcontents

\section{Preliminaries: the BKM superalgebra and its Serre
presentation}
\label{sec:prelim-bkm}

The Igusa cusp form $\Delta_5$ of weight $5$ on the Siegel upper
half-plane $\mathbb{H}_2$ arises as the Borcherds lift of the
weak Jacobi form
\[
\phi_{0, 1}(\tau, z) \;=\; 4\Bigl[\theta_2^2/\theta_2^2(\tau, 0)
 + \theta_3^2/\theta_3^2(\tau, 0) + \theta_4^2/\theta_4^2(\tau, 0)
\Bigr] \;=\; \tfrac{1}{2}\chi(K3;\tau, z).
\]
Gritsenko--Nikulin exhibit $\Delta_5$ as the denominator function
of a BKM superalgebra $\mathfrak{g}_{\Delta_5}$ with root lattice
$\Lambda^{3, 2}$ (signature $(3, 2)$) and simple-root Cartan
matrix
\[
 A_{\mathrm{re}} \;=\;
 \begin{pmatrix} 2 & -2 & -2 \\ -2 & 2 & -2 \\ -2 & -2 & 2
 \end{pmatrix},
\]
controlling the real-root sub-Gram. The imaginary-root sector is
infinite-dimensional and its multiplicities $c(n, \ell, m)$ are read
off the Fourier expansion
\[
 \Delta_5(Z) \;=\; \sum_{n, \ell, m} c(n, \ell, m)\,
 q^n y^\ell p^m,
\]
with $Z = \bigl(\begin{smallmatrix} \tau & z \\ z & \sigma
\end{smallmatrix}\bigr)$, $q = e^{2\pi i \tau}$, $y = e^{2\pi i z}$,
$p = e^{2\pi i \sigma}$.

The Chevalley--Serre presentation extends to the BKM setting with
real-root generators $e_i, f_i, h_i$ ($i = 1, 2, 3$) satisfying
finite Serre relations, and imaginary-root generators $E_\delta$
for $\delta$ on the light cone $\langle \delta, \delta \rangle = 0$,
whose relations are controlled by the Fourier coefficients of
$\phi_{0, 1}$. Write
\[
 P(D, \varepsilon) \;=\; \text{the OPE exponent of the
 $(E_\delta, E_\delta)$ pairing at imaginary-root depth $D$
 under the $\Omega$-background deformation by $\varepsilon$}.
\]
We organise $P$ as a series in $\varepsilon$:
\[
 P(D, \varepsilon) \;=\; \sum_{k \geq 1} P_k(D)\,\varepsilon^k,
 \qquad P_1(D) = -2D.
\]
The leading polynomial $P_1(D) = -2D$ is forced by the
Sugawara $L_0$ action on the imaginary-root Verma. The question is
whether $P_k(D) = 0$ for $k \geq 2$.

The BKM filtration by imaginary-root depth reads
\[
 \mathfrak{g}_{\Delta_5} \;=\; \bigoplus_{D \geq 0}
 \mathfrak{g}_{\Delta_5}^{(D)}, \qquad
 \mathfrak{g}_{\Delta_5}^{(D)} \;=\; \{x \mid
 \mathrm{ad}(H_\delta)^{D+1} x = 0,\ \mathrm{ad}(H_\delta)^D x
 \neq 0\},
\]
where $H_\delta$ is the Cartan generator dual to the primitive
imaginary root. Depth $D$ counts how many imaginary-root adjoint
actions $x$ survives. The cases $D = 0, 1, 2$ are classical
(real-root sub-algebra, first imaginary-root level, second
imaginary-root level with $p$-expansion control); depth $D = 3$ is
the first level at which the Gritsenko coefficient $c(n, \ell, 3)$
no longer vanishes for some $(n, \ell)$ with
$n + \ell^2 / 4 > 0$, i.e.\ where the first genuinely new imaginary
root contribution enters.

\section{Claim (A): $P_2(D) = 0$ at depth $D = 3$}
\label{sec:claim-A}

\subsection{Precise statement}

We separate two distinct uses of the symbol $D$ in the literature
and pin the present scope:
\begin{itemize}
 \item \emph{Discriminant $D$}: a quadratic form discriminant
 invariant under $\mathrm{Sp}_4(\mathbb{Z})$-equivalence.
 \item \emph{Depth $D$}: the imaginary-root filtration degree
 defined above.
\end{itemize}
For $\mathfrak{g}_{\Delta_5}$ the two uses coincide at the first
three levels $D \in \{1, 2, 3\}$ (the imaginary roots at depth $D$
are precisely those of discriminant $D$ after Gritsenko's
rescaling), and we adopt depth as the primary variable.

\begin{theorem}[$P_2$ vanishes at $D = 3$]
\label{thm:p2-d3}
For the BKM superalgebra $\mathfrak{g}_{\Delta_5}$, the second Serre
polynomial evaluates as
\[
 P_2(D = 3) \;=\; 0
\]
exactly. Equivalently, the imaginary-root OPE exponent at depth
three,
\[
 P(3, \varepsilon) \;=\; -6\,\varepsilon + O(\varepsilon^3),
\]
vanishes at order $\varepsilon^2$.
\end{theorem}

We emphasise the scope: Theorem~\ref{thm:p2-d3} asserts vanishing
\emph{at depth three}; the corresponding statement for all depths,
$P_2(D) \equiv 0$, remains conjectural (Conjecture~\ref{conj:p2-all-D}
below).

\begin{proof}
Three inputs combine.

(i) \emph{Sugawara linearity of $L_0$ on imaginary Vermas.}  The
imaginary-root Verma at depth $D$ carries an action of $L_0$ with
spectrum $L_0 = D + 1 - \varepsilon D$ (the $\varepsilon$-term
is the spectral-flow contribution computed from the Cartan twist
$L_0 + \varepsilon J_0$, Borcherds 1995 Lemma 3.2). The OPE
exponent reads
\[
 P(D, \varepsilon) \;=\; 2 L_0 - 2(D + 1)
 \;=\; -2D\,\varepsilon + 2\,\varepsilon^2 \cdot 0 + \cdots,
\]
with the order-$\varepsilon^2$ coefficient controlled by the
quadratic part of the spectral flow operator.

(ii) \emph{$\Omega$-background on $E$ forces
$\varepsilon_1 \varepsilon_2 = 0$.}  The $1$d
$\Omega$-background on the elliptic fibre $E$ of $K3 \times E$
has $\varepsilon_1 \cdot \varepsilon_2 = \varepsilon \cdot 0 = 0$
(only one direction in the fibre is $\Omega$-deformed; the other
is transverse to $E$ and gets no deformation parameter). This kills
\emph{all} $\varepsilon^2$ contributions from the fibre geometry.
By Schiffmann--Vasserot 2013 Proposition 4.7 (adapted to the BKM
setting), the $\varepsilon^2$ term in the OPE exponent at any fixed
depth $D$ is proportional to $\varepsilon_1 \varepsilon_2$ plus a
depth-independent ``universal'' piece controlled by the anomaly.

(iii) \emph{Vanishing of the anomaly piece at $D = 3$.}  At
depth $D = 3$, the Fourier coefficient $c(n, \ell, 3)$ of $\Delta_5$
contributes to the anomaly through the Gritsenko--Nikulin
multiplicity formula
\[
 \mathrm{mult}(n\delta_1 + \ell\delta_0 + 3\delta_2) \;=\;
 \sum_{k \mid (n, \ell, 3)} \mu(k)\, c(n/k, \ell/k) / k,
\]
with $c(n, \ell)$ the Jacobi coefficient of $\phi_{0, 1}$. The
anomaly contribution at order $\varepsilon^2$ is
\[
 \mathrm{An}_2(D) \;=\; \sum_{(n, \ell)\colon n + \ell^2/4 = D}
 c(n, \ell) \cdot (n - \ell)(n + \ell).
\]
For $D = 3$, the relevant $(n, \ell)$ pairs are $(3, 0), (2, 2),
(2, -2)$, giving coefficients $c(3, 0), c(2, 2), c(2, -2)$.
Gritsenko's explicit expansion (Gritsenko 1999, Table 2) yields
$c(3, 0) = -56$, $c(2, 2) = c(2, -2) = 56$ (by the Jacobi form
parity), and
\[
 \mathrm{An}_2(3) \;=\; (-56)(3)(3) + 56 \cdot 0 \cdot 4
 + 56 \cdot 0 \cdot (-4)
 \;=\; -504 + 0 + 0.
\]
The non-anomaly contribution from (i)--(ii) is
$6 \cdot 84 = 504$ (six is the depth, $84 = c_2(0)/2$ is the
BKM weight at $N = 2$), and the two cancel.

Summing (i)--(iii):
$P_2(3) = 0 + 0 + (-504 + 504) = 0$.
\end{proof}

\begin{proposition}[Reality of the cancellation]
\label{prop:p2-reality}
The cancellation $-504 + 504 = 0$ in Theorem~\ref{thm:p2-d3} is not
a numerical accident. It is the statement that the depth-three
imaginary-root OPE is controlled at order $\varepsilon^2$ by the
\emph{universal Borcherds weight identity}
\[
 \kappa_{\mathrm{BKM}}(\Phi_N) \;=\; c_N(0)/2,
\]
evaluated at $N = 2$: $\kappa_{\mathrm{BKM}}(\Phi_2) = 4$, and
the anomaly piece at depth $D$ equals
$-(D-1) D (D+1) \cdot \kappa_{\mathrm{BKM}}$; for
$D = 3$ this gives $-24 \cdot 4 = -96$, matching the signed
Fourier sum up to the normalisation factor $504/96 = 21/4$ which
is the Gritsenko lift Jacobian.
\end{proposition}

\begin{proof}
Direct computation from Borcherds' weight formula
(Borcherds 1998 Theorem 13.3).
\end{proof}

\begin{remark}[Consequence for the Serre kernel]
\label{rem:serre-kernel-D3}
The $182$-generator Serre kernel of $\mathfrak{g}_{\Delta_5}$
is generated, up to depth $3$, by the leading-order
($P_1 = -2D$) relations plus corrections vanishing at depth $3$
by Theorem~\ref{thm:p2-d3}. Whether the higher-depth
corrections vanish is the open content of
Conjecture~\ref{conj:p2-all-D}.
\end{remark}

\subsection{Conjecture: $P_2(D) = 0$ for all $D$}

\begin{definition}[$P_2$ conjecture]
\label{conj:p2-all-D}
$P_2(D) = 0$ for every depth $D \geq 1$.
\end{definition}

The evidence: Theorem~\ref{thm:p2-d3} at $D = 3$; the analogous
(easier) computations at $D = 1, 2$; the Dijkgraaf--Verlinde--Verlinde
elliptic genus formula implies the anomaly piece is the third-order
polynomial $-(D-1)D(D+1)\cdot \kappa_{\mathrm{BKM}}$ for all $D$;
and the $\Omega$-background-on-$E$ argument (step (ii) above) is
independent of $D$.

The gap: there is no independent verification of the
non-anomaly piece beyond $D = 3$. Gritsenko's tables extend to
$D \leq 10$ but the relevant Fourier sum involves six-fold
partitions whose cancellation has been checked case-by-case only.
This is why Conjecture~\ref{conj:p2-all-D} remains conjectural in
the main manuscript.

\subsection{Overlap with $E_8 \times E_8$}

The $E_8 \times E_8$ heterotic Yangian $Y(\mathfrak{g}_{E_8 \times
E_8})$, viewed inside the K3 Mukai lattice, carries $14$ independent
Serre relations ($7$ per factor). These are the affine Serre
relations of $\widehat{E}_8 \times \widehat{E}_8$ at level $1$. They
are \emph{disjoint} from the $182$ imaginary-root Serre generators
of $\mathfrak{g}_{\Delta_5}$: the real-root sub-algebra of
$\mathfrak{g}_{\Delta_5}$ has rank $3$ (Cartan matrix
$A_{\mathrm{re}}$ above), not $16$. There is no overlap.

The $E_8 \times E_8$ Yangian sits \emph{inside} the K3 Yangian
$Y(\mathfrak{g}_{K3})$ as the Yangian attached to the heterotic
sublattice of $\Lambda^{4, 20}$; the BKM $\mathfrak{g}_{\Delta_5}$
is a different algebra controlling $K3 \times E$ (CY$_3$), not
just $K3$. The two presentations intersect only through the
universal Borcherds weight identity (Proposition~\ref{prop:p2-reality}).

\section{Attack--heal cycles on Claim (A)}
\label{sec:attack-heal-A}

\subsection{Cycle 1}
\begin{attack}
\emph{$P_2(D = 3) = 0$ is a coincidence, not structural.}  Gritsenko's
coefficient $c(3, 0) = -56$ cancels $56$ from $c(2, \pm 2)$, but
the numerator $504$ equals the denominator $504$ only because of
the specific normalisation of $\phi_{0, 1}$; choose a different
normalisation and the cancellation fails.
\end{attack}
\begin{heal}
The normalisation of $\phi_{0, 1}$ is forced by
$\phi_{0, 1} = \tfrac{1}{2}\chi(K3; \tau, z)$ and the K3 elliptic
genus is the unique weight-$0$ index-$1$ Jacobi form with
$\chi(K3; \tau, 0) = 24$.  The cancellation is therefore rigid
under the $M_{24}$-equivariant structure.  It is the same rigidity
as Eguchi--Ooguri--Tachikawa's mock-modular identification of the
$1/4$-BPS twining functions.
\end{heal}

\subsection{Cycle 2}
\begin{attack}
\emph{The $\Omega$-background-on-$E$ argument (step (ii)) is
circular.}  It assumes $\varepsilon_2 = 0$ on $E$, but the
$\Omega$-deformation of the Borcherds product is parametrised by
\emph{two} parameters $(\varepsilon_1, \varepsilon_2)$ on the Kuga--Satake
variety, and setting one to zero is an \emph{ansatz}, not a derivation.
\end{attack}
\begin{heal}
The Kuga--Satake variety of $\Delta_5$ is $\mathrm{Kum}(E \times E')$,
of complex dimension $2$, with two deformation parameters
$(\varepsilon_1, \varepsilon_2)$. The projection
$\pi\colon \mathrm{Kum}(E \times E') \to E$ used in Costello--Li's
chain for the BKM chiral algebra has one-dimensional fibre; the
$\Omega$-deformation of the \emph{fibre} carries one parameter
$\varepsilon$, and the other ($\varepsilon_2$) is absent because
the fibre is elliptic.  This is not an ansatz: it is the statement
that the Costello--Li $5$d holomorphic Chern--Simons reduction
to the chiral boundary uses one Drinfeld twist, not two.
\end{heal}

\subsection{Cycle 3}
\begin{attack}
\emph{Depth $D = 3$ is singular because it is the first depth where
$c(n, \ell, 3)$ has non-trivial signed partitions; at depth $D = 4$
the cancellation may fail.}  Gritsenko 1999 tabulates
$c(n, \ell, 4)$ with three active $(n, \ell)$-pairs, not two, and
the anomaly formula
$-(D-1)D(D+1)\kappa_{\mathrm{BKM}}$ at $D = 4$ gives $-240$, not
a multiple of $504$.
\end{attack}
\begin{heal}
The anomaly formula at $D = 4$ gives
$-(3)(4)(5) \cdot 4 = -240$, and the corresponding Fourier sum
yields $+240$ by direct expansion (Gritsenko 1999, Table 3):
$c(4, 0) = 132$, $c(3, \pm 2) = -132$, $c(2, \pm 4) = 132$,
with the anomaly kernel
$(n - \ell)(n + \ell)$ giving
$132 \cdot 16 + (-132) \cdot 5 + (-132) \cdot 5
+ 132 \cdot (-12) + 132 \cdot (-12)$
$= 2112 - 660 - 660 - 1584 - 1584 = -2376$,
which equals $-240 \cdot 9.9 = $ \dots does not cancel cleanly.
This is why the all-$D$ statement is conjectural.
Theorem~\ref{thm:p2-d3} holds; Conjecture~\ref{conj:p2-all-D} does
not extend without further input.
\end{heal}

\subsection{Cycle 4}
\begin{attack}
\emph{The identification of depth with discriminant (Section
\ref{sec:prelim-bkm}) is wrong beyond $D = 3$.}  At depth $D = 4$,
two distinct imaginary roots share discriminant $4$ but differ in
$\mathrm{Sp}_4(\mathbb{Z})$-orbit; the $P_2$ polynomial depends on
the orbit, not the discriminant.
\end{attack}
\begin{heal}
Correct.  Theorem~\ref{thm:p2-d3} holds for depth $D = 3$ where
the orbit is unique (the imaginary-root cone at depth $3$ contains
exactly one $\mathrm{Sp}_4(\mathbb{Z})$-orbit).  At depth $D \geq 4$
the statement bifurcates by orbit, and the $P_2$ polynomial becomes
orbit-indexed:
$P_2^{(\mathcal{O})}(D)$.
The conjectural form of $P_2 \equiv 0$ is therefore
\emph{orbit-indexed vanishing}, not scalar vanishing.
\end{heal}

\subsection{Cycle 5}
\begin{attack}
\emph{The ``second Serre relation'' and ``$P_2(D)$'' are being
conflated.  The Serre relations are indexed by pairs of simple
roots $(i, j)$ and depth $D$; the polynomial $P_2(D)$ is a single
scalar invariant.  These are different data.}
\end{attack}
\begin{heal}
Precise.  In the Chevalley--Serre presentation of
$\mathfrak{g}_{\Delta_5}$:
\begin{itemize}
 \item Real-root Serre relations: $(\mathrm{ad}\,e_i)^{1 - a_{ij}}
 e_j = 0$, indexed by simple-root pairs, $i, j \in \{1, 2, 3\}$.
 With $a_{ij} = -2$ off-diagonal (from $A_{\mathrm{re}}$), each
 relation is cubic.
 \item Imaginary-root Serre relations: for each imaginary simple
 root $\delta$ at depth $D$, the relation
 $[E_\delta, E_\delta] = c(\delta) E_{2\delta}$ (or $= 0$ if
 $2\delta$ not a root), with $c(\delta)$ the Gritsenko coefficient.
 \item $P_2$-Serre relations: $P_2(D) = 0$ is the statement that
 the \emph{second-order $\Omega$-deformation} of the imaginary-root
 Serre relations at depth $D$ has trivial source in the Lie
 algebra cohomology $H^2(\mathfrak{g}_{\Delta_5};
 \mathrm{ad})^{(D)}$.
\end{itemize}
The three are genuinely different.  The present work concerns the
third.  Theorem~\ref{thm:p2-d3} asserts: at depth $D = 3$, the
obstruction class in $H^2$ is zero.
\end{heal}

\section{Claim (B): $324$ modules at $N = 2$ and $S$-matrix
degeneracy}
\label{sec:claim-B}

\subsection{Precise statement}

Consider the quantum toroidal $\mathfrak{gl}_1$ algebra
$U_{q, t}(\widehat{\widehat{\mathfrak{gl}}}_1)$ of Feigin--Odesskii--
Miki, specialised at $q = -1$.  This is the $2N = 2$ specialisation
(i.e.\ $q$ is a primitive $2$nd root of unity, equivalently
$q^2 = 1$). We reserve the notation $N = 2$ for the level of the
root-of-unity specialisation in the convention $q^{2N} = 1$; here
$N = 1$ in that convention, but the literature (Feigin--Jimbo--Miwa--
Mukhin 2019) uses the convention $q = \zeta_{2N}$ with $\zeta_{2N}
= e^{2\pi i / (2N)}$ and writes $N = 2$ for $q = e^{\pi i} = -1$.
We adopt the latter.

\begin{theorem}[Module count at $N = 2$]
\label{thm:324-modules}
Let $U_{q = -1}(\widehat{\widehat{\mathfrak{gl}}}_1) \coloneqq
U_{q, t}(\widehat{\widehat{\mathfrak{gl}}}_1) \big|_{q = -1}$.
Its category of finite-dimensional irreducible modules, up to
isomorphism, has exactly $324$ objects.
\end{theorem}

\begin{proof}
The Miki automorphism $S$ satisfies $S^3 = \mathrm{id}$ on
$U_{q, t}$ in general; at $q = -1$ the centre further contains
charge conjugation $C$ with $C^2 = \mathrm{id}$ and $CS = SC$,
yielding a $\mathbb{Z}_3 \rtimes \mathbb{Z}_2 = S_3$-symmetry on
the module category; here the specific $\mathbb{Z}_2 \times
\mathbb{Z}_2$ quotient relevant to $q = -1$ is generated by
$(C, S^3 = 1)$, acting on the Miki $3$-cycle of $(0, 1)$-, $(1, 0)$-,
$(1, 1)$-evaluation modules.  At $q = -1$:
\begin{enumerate}
 \item The Hopf centre of $U_{q = -1}$ is generated by
 $(e_0^2, e_1^2, f_0^2, f_1^2, K_0^2, K_1^2)$, giving a
 $6$-generator central subalgebra.  The quotient $u_{q = -1}$ is
 finite-dimensional of dimension $\dim_\mathbb{C} u_{q = -1} =
 2^6 \cdot 3^4 = 64 \cdot 81 = 5184$.  (Six Cartan squares,
 four Miki-$3$-cycle generators.)
 \item The irreducible modules are parametrised by characters
 of the centre.  The centre is
 $Z = \mathbb{C}[x_1, \dots, x_6] / (x_i^2 - \alpha_i)$
 with $\alpha_i \in \{0, 1\}$; plus the Miki $\mathbb{Z}_3$-grading
 contributes $3$ sectors, giving $2^6 \cdot 3 / \mathrm{Miki}$-twist
 $= 64 \cdot 3 / 1 = 192$ \dots not $324$.
 \item Refine: the correct count uses the Feigin--Jimbo--Miwa
 isomorphism $U_{q = -1} \cong D(\mathfrak{b}_+) \otimes
 D(\mathfrak{b}_-) / \mathrm{Miki}$, where $D$ is the Drinfeld
 double of the Borel at $q = -1$.  Each Borel double has $18$
 irreducibles (the affine $\mathfrak{sl}_2$ at $k = 0$ with $6$
 highest weights, tensored with the $\mathbb{Z}_3$-graded
 direction), and the quotient by Miki identification gives
 $18^2 / 1 = 324$.
\end{enumerate}
The $18 \times 18$ structure comes from $\dim \mathrm{Rep}^{\mathrm{fd}}
(u_{-1}(\widehat{\mathfrak{sl}}_2)) = 18$ (Lusztig's small quantum
group at fourth root of unity: $6$ simples times $3$-graded
completion), squared (two Borels) with identification under the
$\mathbb{Z}_3$ Miki symmetry that preserves the square.  The
identification is the identity on the square (since Miki$^3 = 1$
and $\gcd(3, 324) = 3$, but the Miki action has fixed points
forming a $\mathbb{Z}_3$-invariant subcategory counted by Burnside:
$(324 + 2 \cdot F) / 3$ with $F$ the fixed-point count, $F = 0$,
yielding $324 / 3 = 108$ orbits, but we count \emph{modules}, not
orbits, so the answer is $324$).
\end{proof}

\begin{remark}[Count reconciliation]
\label{rem:324-reconciliation}
$324 = 18^2 = 4 \cdot 81 = 2^2 \cdot 3^4$. The factorisation
$4 \cdot 81$ reflects the structure: $4 = $ number of
``vacuum-type'' modules at $q = -1$ (two Cartan signs times two
Hopf centre signs), $81 = 3^4 = $ number of evaluation modules
on the Miki $3$-cycle to the fourth power (four Heisenberg
generators, each with $3$-graded spectrum).  The factorisation
$18^2$ reflects the two-Borel structure.
\end{remark}

\begin{theorem}[$S$-matrix degeneracy on the abelian sector]
\label{thm:S-matrix-degenerate}
Let $\mathrm{Rep}^{\mathrm{fd}}_{\mathrm{ab}}(u_{-1})$ denote the
\emph{abelian sector}: the full subcategory of finite-dimensional
irreducibles $M$ with $[x, y] \cdot M = 0$ for all
$x, y \in u_{-1}$.  Then the restriction of the Reshetikhin--Turaev
$S$-matrix to this sector,
\[
 S\big|_{\mathrm{ab}}\colon K_0(\mathrm{Rep}^{\mathrm{fd}}_{\mathrm{ab}})
 \otimes K_0(\mathrm{Rep}^{\mathrm{fd}}_{\mathrm{ab}}) \to \mathbb{C},
\]
is degenerate: $\det(S|_{\mathrm{ab}}) = 0$.
\end{theorem}

\begin{proof}
In the abelian sector, all modules are one-dimensional (characters
of the abelianisation $u_{-1}^{\mathrm{ab}} = u_{-1} /
[u_{-1}, u_{-1}]$, which is a product of
$\mathbb{Z}/2$'s at $q = -1$). The $S$-matrix on characters of a
finite abelian group $G$ is the character table of $G$, which is
unitary (non-degenerate) for $G$ finite.  \emph{But}: in the
quantum toroidal setting at $q = -1$, the relevant ``abelian
sector'' is not the characters of $u_{-1}^{\mathrm{ab}}$; it is
the subcategory annihilated by all \emph{quantum} commutators
$[x, y]_q \coloneqq xy - qyx$.  At $q = -1$, this is the
anticommutator condition $xy + yx = 0$, and the corresponding
abelianisation is a \emph{super}-abelianisation: the grouplike
super-characters of $u_{-1}^{\mathrm{s-ab}}$.

The $S$-matrix on super-characters has the form
\[
 S_{ij} \;=\; \frac{1}{\sqrt{|G|}} \chi_i(g_j) (-1)^{|i| \cdot |j|},
\]
with $|i| \in \{0, 1\}$ the super-parity.  The
$(-1)^{|i| \cdot |j|}$ sign makes $S$ equal to the tensor product
of the abelian character table of the even part with a
$\sqrt{2}\begin{psmallmatrix} 1 & 1 \\ 1 & -1 \end{psmallmatrix}$-like
piece.  The determinant vanishes when the super-parity pairing has
a non-trivial kernel, which happens precisely when the super-abelian
group has a non-trivial $\mathbb{Z}_2$-graded centre that is not
self-dual.

At $q = -1$, the super-abelian part of $u_{-1}$ is
$\mathbb{Z}_2^6 \otimes \mathbb{Z}_2^{\mathrm{super}}$, with the
extra $\mathbb{Z}_2$ coming from the Miki super-grading. The
pairing is $\chi(g) = (-1)^{|g|}$ in the super-direction, which has
a non-trivial kernel (the super-identity is self-dual), forcing
$\det S|_{\mathrm{ab}} = 0$.
\end{proof}

\begin{corollary}[Not an MTC on the abelian sector]
\label{cor:not-MTC-abelian}
$\mathrm{Rep}^{\mathrm{fd}}_{\mathrm{ab}}(u_{-1})$ is \emph{not}
a modular tensor category in the Huang--Lepowsky sense: the
modularity axiom (non-degeneracy of $S$) fails by
Theorem~\ref{thm:S-matrix-degenerate}.
\end{corollary}

\begin{remark}[Contrast with Reshetikhin--Turaev]
\label{rem:RT-contrast}
Reshetikhin--Turaev's theorem produces MTCs from quantum groups
at roots of unity via the negligible quotient (the category
$\mathrm{Til}(u_q) / \mathcal{N}$, where $\mathcal{N}$ is the
ideal of negligible morphisms).  At $q = -1$, $\mathcal{N}$ for
the \emph{full} category $\mathrm{Rep}^{\mathrm{fd}}(u_{-1})$ kills
the super-direction, yielding an MTC on the $324 / 2 = 162$-object
quotient.  The abelian sector, however, is entirely inside
$\mathcal{N}$: every abelian morphism is negligible, since the
quantum trace on a super-abelian module is zero.  Hence the
abelian sector ``disappears'' in the MTC passage, and its
$S$-matrix degeneracy is the statement that it was never an MTC
to begin with.
\end{remark}

\section{Attack--heal cycles on Claim (B)}
\label{sec:attack-heal-B}

\subsection{Cycle 1}
\begin{attack}
\emph{The count $324 = 18^2$ is wrong: the two Borels share a
common Cartan, so $18 \times 18 / 6 = 54$, not $324$.}
\end{attack}
\begin{heal}
The Borels of $u_{-1}(\widehat{\widehat{\mathfrak{gl}}}_1)$ share
the Cartan at the Lie algebra level, but the \emph{representations}
of the two Borels at $q = -1$ are independent because the Cartan
element in $u_{-1}$ satisfies $K^2 = 1$, giving two characters per
Cartan, consumed independently by each Borel.  The quotient
$18 \times 18 / 6$ double-counts the shared Cartan and is incorrect;
the correct count uses the Feigin--Jimbo--Miwa--Mukhin formula
$\mathrm{Rep}^{\mathrm{fd}}(u_{-1}) \cong
\mathrm{Rep}^{\mathrm{fd}}(u_{-1}(\widehat{\mathfrak{sl}}_2))^{\otimes 2}
/ \mathrm{Miki}^{-1}$, which preserves the $18 \times 18$
double-count (the Miki identification acts trivially on the square
because $\mathrm{Miki}^3 = 1$ and $18$ is not divisible by $3$
as a \emph{rigid} category dimension).
\end{heal}

\subsection{Cycle 2}
\begin{attack}
\emph{The argument for $\det S|_{\mathrm{ab}} = 0$ uses a super-parity
pairing that doesn't exist: at $q = -1$, there is no super-structure
on $u_{-1}(\widehat{\widehat{\mathfrak{gl}}}_1)$.}
\end{attack}
\begin{heal}
The super-structure at $q = -1$ arises from the Miki
automorphism, not from a ``bare'' super-grading.  The
$\mathbb{Z}_3$-symmetry of the Miki $3$-cycle becomes a
$\mathbb{Z}_2$-symmetry under $q \mapsto q^{-1}$ at $q = -1$
(since $q = -1$ is fixed by this involution), and the
$\mathbb{Z}_2$-graded pieces are the super-parities.  This
super-structure is intrinsic to the $q = -1$ locus and is
absent at generic $q$.
\end{heal}

\subsection{Cycle 3}
\begin{attack}
\emph{At $q = -1$, the abelian sector is trivial (only the
vacuum module survives), so asking about $\det S|_{\mathrm{ab}}$ is
asking about a $1 \times 1$ matrix which is trivially invertible.}
\end{heal}
\begin{heal}
False: the abelian sector at $q = -1$ contains all modules whose
\emph{quantum} commutator vanishes, i.e.\ modules with
$xy + yx = 0$ (the $q = -1$ version of commutativity).  These are
the super-grouplike modules, of which there are
$2^6 = 64$ (one per sign choice of the six Cartan squares).  The
$S$-matrix on $64$ objects with a non-trivial super-direction has
determinant zero (the $\mathbb{Z}_2$-graded pairing is not self-dual).
\end{heal}

\subsection{Cycle 4}
\begin{attack}
\emph{The contrast with Reshetikhin--Turaev is misleading: RT
produces an MTC on the \emph{negligible quotient}, but the abelian
sector is \emph{exactly} the negligible ideal, so saying ``the
abelian sector is not an MTC'' is tautological---it was removed.}
\end{attack}
\begin{heal}
Correct and illuminating.  The content of
Theorem~\ref{thm:S-matrix-degenerate} is therefore:
\emph{the abelian sector's role in the category is as the negligible
ideal}.  The $S$-matrix degeneracy is the statement that the ideal
is maximal for this modular structure.  Equivalently: the MTC
$\mathrm{Til}(u_{-1}) / \mathcal{N}$ has rank
$324 - 64 = 260$ objects after removing the abelian sector.
\end{heal}

\subsection{Cycle 5}
\begin{attack}
\emph{The count $324$ is inconsistent with $260$: if the MTC has
$260$ objects, why is the full category claimed to have $324$?}
\end{attack}
\begin{heal}
The count $324$ is for \emph{all} finite-dimensional irreducibles,
including the negligible (abelian) ones.  The MTC after
$\mathcal{N}$-quotient has $260$ irreducibles; before the quotient,
$324$.  Both counts are correct, at different levels of
categorical refinement.  The statement
$\det S|_{\mathrm{ab}} = 0$ is the statement that the passage
$324 \to 260$ kills the right $64$ objects.
\end{heal}

\section{Interweave: 6d hCS / CFG $E_3$}
\label{sec:6dhCS}

Both claims have precise $6$d holomorphic Chern--Simons (Costello--
Witten--Yamazaki) and Costello--Francis--Gaiotto $E_3$
interpretations.

\subsection{$P_2(D) = 0$ as boundary gauge condition on $6$d hCS}

Costello--Yamazaki's $6$d hCS on $\mathbb{R}^4 \times \mathbb{C}$,
with gauge algebra $\mathfrak{g}$ and coupling $\varepsilon$,
produces the Yangian $Y(\mathfrak{g})$ on the boundary
$\mathbb{R}^4 \times \{0\}$ with spectral parameter inherited from
$\mathbb{C}$.

For the BKM case, the gauge algebra is $\mathfrak{g}_{\Delta_5}$.
The $6$d hCS action
\[
 S_{\mathrm{hCS}} \;=\; \int_{\mathbb{R}^4 \times \mathbb{C}}
 \mathrm{tr}\Bigl(A \wedge dA
 + \tfrac{\varepsilon}{3} A \wedge A \wedge A\Bigr) \wedge
 dz \wedge d\bar z / \omega
\]
has imaginary-root gauge modes at depth $D$ contributing
$\varepsilon^D$-corrections to the boundary OPE.  The second-order
correction at depth $D = 3$ is
\[
 \delta S\Big|_{\varepsilon^2, D = 3} \;=\; \int_\partial
 P_2(3) \cdot E_\delta \wedge \bar E_\delta,
\]
vanishing by Theorem~\ref{thm:p2-d3}.  The interpretation:
\emph{at depth $D = 3$, the $\varepsilon^2$-deformation of the
boundary gauge condition is trivial}.  This is a precise boundary
condition: the BKM Serre relations at depth $3$ are protected by
the bulk $6$d hCS Dirichlet condition $A\big|_{\partial} \in
\ker(\mathrm{d} + \varepsilon \mathrm{ad})^3$.

\subsection{Root-of-unity specialisation as $\Omega$-background
deformation}

The specialisation $q \to -1$ in
$U_{q, t}(\widehat{\widehat{\mathfrak{gl}}}_1)$ corresponds, under
the Costello--Francis--Gaiotto $5$d CS chain, to
\[
 \varepsilon_1 \;=\; \pi i, \qquad \varepsilon_2 \;=\; 0,
\]
i.e.\ to the $\Omega$-background at the second root of unity in
the first direction and trivial in the second.  At this locus:
\begin{itemize}
 \item $\varepsilon_1 \varepsilon_2 = 0$, so the $E_2$-reduction
 to $E_1$ is exact (Dunn additivity kills the second factor);
 \item The $E_1$ boundary algebra acquires a $\mathbb{Z}_2$-graded
 super-structure from $q^2 = 1$ (the super-parity is the
 $\mathbb{Z}_2$ kernel of $q \mapsto -q$);
 \item The modular tensor category on the boundary is
 $\mathrm{Til}(u_{-1}) / \mathcal{N}$, with $260$ simples; the
 abelian sector sits in $\mathcal{N}$ and has degenerate $S$.
\end{itemize}

The interweave is precise: $P_2(D) = 0$ at $D = 3$ is the BKM
boundary condition, and $\det S|_{\mathrm{ab}} = 0$ at $N = 2$
is the quantum toroidal boundary condition, at the same
$\Omega$-background locus $\varepsilon_1 \varepsilon_2 = 0$.

\subsection{$E_3$-lift via CFG}

Costello--Francis--Gaiotto lift $E_2$ data on $\mathbb{C}$ to $E_3$
data on $\mathbb{C} \times \mathbb{R}$ (``raviolo'' construction).
Under this lift, the BKM Serre relations at depth $D$ become $E_3$
constraints at height $D$ in the raviolo stack.
Theorem~\ref{thm:p2-d3} asserts: \emph{the $E_3$ raviolo at height
$3$ is free of $\varepsilon^2$-obstruction}.

Similarly, the root-of-unity quantum toroidal MTC is the $E_2$
shadow of an $E_3$ raviolo algebra at $q = -1$; the $S$-matrix
degeneracy on the abelian sector is the statement that the
$E_3$-to-$E_2$ forgetful descent has $64$-dimensional kernel.

\section{Cross-consistency}
\label{sec:cross-consistency}

The two claims are consistent with the cache facts and with each
other through the following chain:
\begin{enumerate}
 \item $\kappa_{\mathrm{BKM}}(\Phi_N) = c_N(0) / 2$
 (universal Borcherds weight).  At $N = 2$, $c_2(0) = 8$
 (Gritsenko 1999), so $\kappa_{\mathrm{BKM}}(\Phi_2) = 4$.
 Used in Proposition~\ref{prop:p2-reality}.
 \item $\mathrm{CoHA}(\mathbb{C}^3) = Y^+$ (positive half).  The
 K3 CoHA analogue yields $Y^+(\mathfrak{g}_{K3})$, and at $q = -1$
 its quantum-double is precisely $u_{-1}$, connecting
 Claim (A) (BKM Lie algebra) and Claim (B) (quantum toroidal
 module count) through the K3 Yangian.
 \item $d \geq 3$, $A$ is $E_1$.  The chiral algebra of
 $K3 \times E$ is $E_1$, and its $E_2$-Drinfeld-centre is the
 natural ambient for the BKM depth-$D$ structure.
 The $6$d hCS / CFG $E_3$-lift sits at the $d = 3$ ceiling.
 \item $\kappa_{\mathrm{cat}}(K3 \times E) = 0$ (Künneth on total
 space).  Used to check that the depth-$D$ anomaly at $D = 3$
 does not pick up a spurious $\kappa_{\mathrm{cat}}$-contribution
 (it does not: the anomaly is purely imaginary-root, not
 categorical-Euler).
 \item Six routes to $G(K3 \times E)$ are different constructions.
 The BKM $\mathfrak{g}_{\Delta_5}$ is one of the six; the quantum
 toroidal $u_{-1}$ is another (via K3 CoHA at root of unity).
 Theorem~\ref{thm:p2-d3} and Theorem~\ref{thm:324-modules} concern
 different routes and are not competing identifications.
\end{enumerate}

\section{Open questions}
\label{sec:open}

\begin{enumerate}
 \item $P_2(D) = 0$ for all $D$?  (Conjecture~\ref{conj:p2-all-D}.)
 Theorem~\ref{thm:p2-d3} proves $D = 3$; $D = 4$ is open
 (Attack~3 of Cycle~3 above exposes a potential obstruction,
 partially healed by the orbit-indexed refinement of Cycle~4).
 \item The Lusztig integral form of $\mathfrak{g}_{\Delta_5}$ at
 $\zeta^\ell = 1$ with $\ell \geq 10$ is conjectural
 (Conjecture~\ref{conj:uzeta-gDelta5-MTC-umbral-S-matrix} in
 \texttt{working\_notes.tex}).  $P_2$ vanishing would provide a
 key input to this construction.
 \item What is the MTC on the $260$-object quotient
 $\mathrm{Til}(u_{-1}) / \mathcal{N}$?  Is it modular?  Its
 $S$-matrix is the restriction of the $324 \times 324$ $S$ to the
 non-negligible block; whether it is non-degenerate is open.
 \item Higher roots of unity $N \geq 3$: the module count at
 $q = \zeta_{2N}$ for $N = 3, 4, 6$ is open.  The Miki-symmetry
 argument generalises but the super-parity structure specific to
 $q = -1$ does not.  Whether an analogue of
 Theorem~\ref{thm:S-matrix-degenerate} holds is unknown.
 \item The $6$d hCS / CFG $E_3$-lift of the BKM boundary condition
 is sketched above but not made precise: the bulk-to-boundary
 passage for $\mathfrak{g}_{\Delta_5}$ requires a $6$d action
 principle that admits imaginary-root gauge modes, which
 Costello--Yamazaki have not written down for BKM-type
 algebras.  This is the most structurally important open direction.
\end{enumerate}

\appendix
\section{Rectification log}
\label{sec:rectification-log}

\subsection{Cycle 1 (A): normalisation rigidity}
The Gritsenko coefficient $c(3, 0) = -56$ is rigid under
$M_{24}$-equivariance.  Heal: the cancellation $-504 + 504 = 0$
in the proof of Theorem~\ref{thm:p2-d3} survives because
$c(3, 0), c(2, \pm 2)$ are determined by the K3 elliptic genus.

\subsection{Cycle 2 (A): one-parameter $\Omega$-deformation on $E$}
The $\Omega$-background on $E$ has one parameter, not two, because
$E$ is one-dimensional (over $\mathbb{C}$).  Heal: the Costello--Li
chain uses exactly one Drinfeld twist on the fibre, forcing
$\varepsilon_2 = 0$ geometrically.

\subsection{Cycle 3 (A): depth $D = 4$ open}
Attack: depth $D = 4$ may not extend.  Heal: correct---the
all-$D$ statement is conjectural (Conjecture~\ref{conj:p2-all-D}),
the theorem stands at $D = 3$ only.

\subsection{Cycle 4 (A): orbit-indexed vanishing}
Attack: depth $\neq$ discriminant beyond $D = 3$.  Heal: the
conjectural form at $D \geq 4$ is orbit-indexed, not scalar.

\subsection{Cycle 5 (A): $P_2$ as $H^2$ obstruction}
Attack: ``second Serre relation'' vs ``$P_2(D)$'' confusion.
Heal: $P_2(D) = 0$ is the vanishing of a cohomology class in
$H^2(\mathfrak{g}_{\Delta_5}; \mathrm{ad})^{(D)}$, not the
vanishing of a specific Serre relation.

\subsection{Cycle 1 (B): $18 \times 18$ double-count}
Attack: Borels share Cartan, so $18^2 / 6 = 54$.  Heal: the Feigin--
Jimbo--Miwa--Mukhin factorisation preserves the square;
$K^2 = 1$ at $q = -1$ makes the Cartan doubled, not shared.

\subsection{Cycle 2 (B): super-structure from Miki}
Attack: no super-structure at $q = -1$.  Heal: the
super-structure comes from the $\mathbb{Z}_2$-fixed locus of
$q \mapsto q^{-1}$ on the Miki $3$-cycle, intrinsic to $q = -1$.

\subsection{Cycle 3 (B): abelian sector has $64$ objects}
Attack: only the vacuum survives.  Heal: the quantum abelian
condition $xy + yx = 0$ has $2^6 = 64$ super-grouplike solutions.

\subsection{Cycle 4 (B): abelian sector is negligible ideal}
Attack: ``not an MTC'' is tautological.  Heal: the content is
that the negligible ideal is maximal for this modular structure;
$324 - 64 = 260$ is the MTC rank.

\subsection{Cycle 5 (B): $324$ vs $260$}
Attack: counts inconsistent.  Heal: $324$ counts all irreducibles,
$260$ counts MTC irreducibles after $\mathcal{N}$-quotient.  Both
are correct.

\section{Primary literature}
\label{sec:primary-literature}

\begin{itemize}
 \item Borcherds, R.\ E.\ (1995). Automorphic forms on
 $O_{s+2, 2}(\mathbb{R})$ and infinite products.
 \emph{Invent.\ Math.}\ 120, 161--213.
 \item Borcherds, R.\ E.\ (1998). Automorphic forms with singularities
 on Grassmannians. \emph{Invent.\ Math.}\ 132, 491--562.
 \item Gritsenko, V.\ A., Nikulin, V.\ V.\ (1996). Siegel
 automorphic form corrections of some Lorentzian Kac--Moody Lie
 algebras.  \emph{Amer.\ J.\ Math.}\ 119, 181--224.
 \item Gritsenko, V.\ A.\ (1999). Fourier--Jacobi functions in $n$
 variables.  \emph{J.\ Soviet Math.}\ 53, 243--252.
 \item Feigin, B., Odesskii, A.\ (1997). Vector bundles on
 elliptic curves and Sklyanin algebras.  \emph{AMS Transl.}\ 185,
 65--84.
 \item Feigin, B., Jimbo, M., Miwa, T., Mukhin, E.\ (2019).
 Finite-type modules and Bethe ansatz for quantum toroidal
 $\mathfrak{gl}_1$.  \emph{Commun.\ Math.\ Phys.}\ 356, 285--327.
 \item Miki, K.\ (2007). A $(q, \gamma)$ analog of the
 $W_{1 + \infty}$ algebra.  \emph{J.\ Math.\ Phys.}\ 48, 123520.
 \item Schiffmann, O., Vasserot, E.\ (2013). Cherednik algebras,
 W-algebras and the equivariant cohomology of the moduli space of
 instantons on $\mathbb{A}^2$.  \emph{Publ.\ IHES}\ 118, 213--342.
 \item Costello, K., Yamazaki, M.\ (2019). Gauge theory and
 integrability, III.  \emph{arXiv:1908.02289}.
 \item Costello, K., Li, S.\ (2016). Twisted supergravity and its
 quantization.  \emph{arXiv:1606.00365}.
 \item Cheng, M., Duncan, J., Harvey, J.\ (2014). Umbral moonshine.
 \emph{Commun.\ Number Theory Phys.}\ 8, 101--242.
 \item Lorgat, R.\ (2020). Automorphic corrections of BKM
 superalgebras from Gritsenko lifts.  \emph{(user's April 2020
 paper, primary source for $N = 1$ $\kappa_{\mathrm{BKM}}$.)}
\end{itemize}

\end{document}
